\begin{frame}{Funções de transferência e polos}
  \small
  \begin{columns}[T,onlytextwidth]
    \begin{column}{0.48\textwidth}
      \begin{block}{Circuito 1}
        \[
          H_1(s)=\frac{-470000}
          {10{,}3823s^2+2650{,}8s+47000}
        \]
        \[
          p_1\approx-19{,}17\,\mathrm{s}^{-1},\qquad
          p_2\approx-236{,}15\,\mathrm{s}^{-1}.
        \]
        \begin{itemize}
          \item $K_1=-10$; magnitude em DC: $20\,\mathrm{dB}$.
          \item $\omega_{n1}\approx67{,}28\,\mathrm{rad/s}$.
          \item $\zeta_1\approx1{,}897$: sobreamortecido.
          \item Resposta ao degrau monótona.
        \end{itemize}
      \end{block}
    \end{column}
    \begin{column}{0.48\textwidth}
      \begin{block}{Circuito 2}
        \[
          H_2(s)=\frac{-470000}
          {10{,}3823s^2+2650{,}8s+470000}
        \]
        \[
          p_{1,2}\approx-127{,}66\pm j170{,}21\,\mathrm{s}^{-1}.
        \]
        \begin{itemize}
          \item $K_2=-1$; magnitude em DC: $0\,\mathrm{dB}$.
          \item $\omega_{n2}\approx212{,}77\,\mathrm{rad/s}$.
          \item $\zeta_2=0{,}6$: subamortecido.
          \item Oscilação amortecida e sobressinal.
        \end{itemize}
      \end{block}
    \end{column}
  \end{columns}

  \vspace{0.12cm}
  \centering\scriptsize
  Polos obtidos de $\operatorname{den}H(s)=0$.
  As partes reais negativas garantem estabilidade e o decaimento dos termos transitórios.
\end{frame}
